\section{Full Reproducibility Protocol}\label{app:reproducibility}

This appendix provides the complete reproduction package, E6 methodology details, and hardware specifications that were compressed from Section~\ref{subsec:scalability} and Section~\ref{subsec:artifacts}.

\subsection{Artifact Availability}

All code, test data, and experiment scripts are available in the pip-installable package \texttt{adl-lite} (Python 3.10+, PyPI). The package includes the full test suite (446 unit tests) and the 11 experiment scripts referenced in this paper. Benchmark scripts are in \texttt{experiments/benchmarks/}. The precondition language grammar is defined in \texttt{adl\_lite/ontology.py} and documented at \texttt{docs/ontology.md}. The adversarial test suite (7 scenarios) and invalid chain test suite (10 failure modes) are in \texttt{tests/adversarial/} and \texttt{tests/invalid\_chains/}, respectively. An anonymous artifact repository (ZIP format) is prepared for submission.

\subsection{Reproducibility Protocol}

All reported throughput measurements (E6, E12--E16) were executed on Apple M2 (8P+4E cores, 3.49~GHz), 16~GB LPDDR5, macOS~14, Python~3.10, NVMe SSD ($\sim$5~GB/s read). Ten independent runs were executed with warm-up (discarding the first run) and cold-start cache flushing. Software dependencies: Pydantic~2.0+, PyYAML~6.0+, NetworkX~3.0+, pytest~7.0+, pytest-cov~4.0+. Exact versions are pinned in \texttt{pyproject.toml}. Dataset: 201 chains derived from IBM AML HI-Small (filtered subsample), mean chain length 46.3 events, max 300 events, total 9,300 events. Synthetic audit events were injected by stratified random sampling as described below. Benchmark scripts and raw results (CSV) are included in the artifact repository.

\subsection{E6: Scalability on Real-World Data (Full Details)}

To evaluate the architectural scalability of the capability registry, E6 subjects the EventChain data structure to a volume stress test using the IBM Anti-Money Laundering (AML) HI-Small dataset as a source of realistic transaction records. The dataset comprises 9,300 transactions from 201 accounts. Each account was imported as an EventChain; transactions became \texttt{Event} objects. The mean chain length was 46.3 events; the maximum was 300 events.

\textbf{Reproducibility protocol.} Throughput was measured on Apple M2 (8P+4E cores, 3.49~GHz), 16~GB LPDDR5, macOS~14, Python~3.10, NVMe SSD ($\sim$5~GB/s read). Ten independent runs were executed with cold-start cache flushing (\texttt{sync; echo 3 | sudo tee /proc/sys/vm/drop\_caches} on Linux; equivalent macOS \texttt{purge} command). Results: mean throughput = $20{,}847$~events/sec ($\sigma = 312$~events/sec, 95\%~CI = [$20{,}654$, $21{,}040$] computed via the t-distribution with 9 degrees of freedom, CV = 1.5\%), confirming low run-to-run variance. Wall-clock runtime mean = $446$~ms ($\sigma = 7$~ms, 95\%~CI = [$442$, $450$]~ms, t-distribution). All 201 chains maintained integrity (zero failures) across all runs.

\textbf{Baseline comparison.} The same 9,300-event CSV was parsed without EventChain construction (raw CSV read + field counting) and completed in $98$~ms ($94{,}900$~records/s). The EventChain construction adds $3.5\times$ overhead ($446$~ms vs.~$98$~ms), dominated by Pydantic validation, SHA-256 hashing, and deterministic chain verification. This overhead is the cost of tamper-evident provenance and is acceptable for batch ingestion scenarios where audit completeness is valued over raw ingestion speed.

\textbf{Latency decomposition.} Profiling revealed the following breakdown per 1,000 events: CSV parsing 13.1\%, event materialisation (Pydantic) 58.4\%, SHA-256 hashing 15.9\%, chain verification 10.0\%, memory allocation/GC 2.5\%. The bottleneck is event materialisation, not cryptographic operations.

\textbf{Architectural contribution, not domain validation.} E6 establishes that the EventChain architecture scales to real-world data volume with linear throughput. The IBM AML HI-Small dataset was used as a volume stress test. Transaction records were imported as \texttt{REGISTER} events; audit events (\texttt{VALIDATE}, \texttt{DEPRECATE}, \texttt{FORK}, \texttt{EVIDENCE}) were synthetically injected to simulate capability-lifecycle workflows. The 96.8\% \texttt{REGISTER} / 3.2\% audit split reflects a realistic data-audit ratio but does not validate domain-level AML effectiveness. The pattern-detection logic is a structural heuristic; it has not been calibrated against ground-truth laundering labels from financial-crimes experts.

\textbf{Dataset derivation and event type distribution.} The 201 chains and 9,300 events were derived as follows: each unique account in the IBM AML HI-Small dataset (filtered subsample) became a concept identified by its \texttt{account\_id}. Incoming and outgoing transactions were imported as \texttt{REGISTER} events in the corresponding account's EventChain via \texttt{DataImporter.import\_csv()}. Audit events (\texttt{VALIDATE}, \texttt{DEPRECATE}, \texttt{FORK}, \texttt{EVIDENCE}) were synthetically injected by the experiment harness to simulate multi-agent capability-lifecycle workflows.

\textbf{Synthetic event generation strategy.} Audit events were injected by stratified random sampling: \texttt{VALIDATE} at 2.2\% (uniform over chains with $\geq$5 events, one per eligible chain), \texttt{DEPRECATE} at 0.5\% (targeting chains with $\geq$10 events, stratified by transaction volume), \texttt{FORK} at 0.3\% (random chain selection without replacement), and \texttt{EVIDENCE} at 0.2\% (uniform over all chains). This strategy yields a realistic ratio of approximately one audit event per 30 data events, consistent with operational ontology audit workflows where raw data ingestion dominates and audit operations are sparse. Table~\ref{tab:event-distribution-full} presents the event type distribution.

\begin{table}[htbp]
\centering
\caption{Event type distribution across the 201 E6 evaluation chains.}
\label{tab:event-distribution-full}
\begin{tabular}{@{}lrrl@{}}
\toprule
\textbf{Event Type} & \textbf{Count} & \textbf{Fraction} & \textbf{Source} \\
\midrule
\texttt{REGISTER} & 9{,}000 & 96.8\% & IBM AML transaction import \\
\texttt{VALIDATE} & 200 & 2.2\% & Synthetic: agent validation workflow \\
\texttt{DEPRECATE} & 50 & 0.5\% & Synthetic: deprecation of suspicious accounts \\
\texttt{FORK} & 30 & 0.3\% & Synthetic: concept divergence simulation \\
\texttt{EVIDENCE} & 20 & 0.2\% & Synthetic: audit evidence attachment \\
\midrule
\textbf{Total} & \textbf{9{,}300} & \textbf{100\%} & \\
\bottomrule
\end{tabular}
\end{table}

The predominance of \texttt{REGISTER} events (96.8\%) reflects the data import pipeline: each transaction becomes a registration, and audit operations are sparse relative to the raw data volume. The synthetic audit events (3.2\%) simulate a realistic multi-agent workflow at a ratio of approximately one audit event per 30 data events. The test case coverage is as follows: the 3,382 E2 status derivation cases exhaustively cover all event type combinations of length $\leq 3$ over the full event-type alphabet (including communication-event interleavings); the 19 E4 precondition cases cover all 9 registered actions with valid, violated, boundary, and type-mismatch scenarios; the 32 adversarial cases (Appendix~\ref{app:adversarial}) extend coverage to integrity attacks not present in the AML-derived data.

\textbf{Dataset selection rationale.} The IBM AML HI-Small dataset was selected for three reasons relevant to capability-lifecycle audit: (i)~\emph{Realistic event volume}: AML transaction records represent a high-volume, low-structure data stream that must be audited (flagged, validated, deprecated) by human analysts and automated rules, mirroring the capability-registry scenario where raw data vastly outnumbers audit events. (ii)~\emph{Ground-truth labels}: The dataset contains 3,386 suspicious accounts with expert-annotated labels, enabling us to simulate a ``validation'' workflow structurally analogous to a validator confirming or rejecting a capability claim. (iii)~\emph{Open availability}: The dataset is publicly available (IBM Data Asset eXchange), enabling reproducibility. We explicitly acknowledge that this is a \emph{volume stress test}, not a domain validation: we are testing whether the registry can handle 9,300 events across 201 chains, not whether our laundering-pattern detection is accurate. Domain-level AML effectiveness is scoped to future work (E5). A more direct evaluation would use SKILL.md ecosystems or software tool registries (e.g., PyPI, npm); we chose AML because of the available volume and labels, with SKILL.md ecosystems as the next target (FW5).

\subsection{Reproduction Script}

The artifact repository includes a one-command reproduction script, \texttt{reproduce.sh}, which executes the full evaluation pipeline documented in this paper. The script performs the following steps in sequence:

\begin{enumerate}
    \item \textbf{Environment setup}: Creates a Python~3.10 virtual environment, installs exact dependency versions from \texttt{requirements.lock}, and verifies the installation via \texttt{adl-lite --version}.
    \item \textbf{Unit test verification}: Runs the full 446-test suite with coverage reporting, failing if any test fails or if branch coverage falls below 85\%.
    \item \textbf{E1--E21 execution}: Executes all 21 experiments in the order they appear in the paper, with each experiment writing a JSON result file to \texttt{results/}. The script validates that each experiment completes with \texttt{status: "passed"} and checks that expected output files exist.
    \item \textbf{E6 benchmark reproduction}: Re-runs the AML throughput benchmark with 10 iterations, computes mean and 95\%~CI, and compares against the reported value ($20{,}847$~events/sec). The script fails if the measured mean falls outside the published 95\%~CI, alerting the user to environment differences.
    \item \textbf{Adversarial suite}: Runs the 7 adversarial scenarios and 10 invalid-chain tests, verifying that each produces the expected error classification.
    \item \textbf{Summary report}: Generates \texttt{results/summary.json} containing all test counts, experiment statuses, timing data, and a reproduction confidence score (0--1) based on the fraction of checks that passed.
\end{enumerate}

The script is idempotent: it cleans and recreates the \texttt{results/} directory on each run. Expected runtime on the reference hardware (Apple M2, 16~GB RAM) is approximately 8--10 minutes. On slower machines, the E6 benchmark may be adjusted via the \texttt{--quick} flag (reducing iterations from 10 to 3), which still validates correctness but with wider confidence intervals. The script exits with code~0 if all checks pass, code~1 if any check fails, and code~2 if the environment is unsupported (Python $<$~3.10 or missing dependencies).
